\section{Training dynamics and loss-landscape geometry}
\label{sec:dynamics}

High-dimensional theory is at its strongest for equilibrium properties of a
trained network. It has much less to say about the training \emph{process}:
when generalisation arrives relative to memorisation, why test risk can be
non-monotone in model size, whether a network stays near its linearised
trajectory or reshapes its features, whether independently trained solutions
share a basin, and whether minibatch noise acts as a genuine temperature. Each
family below turns one of those questions into a single measured number.

\subsection{The lazy--rich crossover}

\manuscriptfigure{figures/family_lazy_rich.pdf}{%
  \textbf{Feature learning as a function of initialisation scale.} The order
  parameter is the centred alignment \eqref{eq:lazy} between the representation
  kernel at initialisation and after training: alignment near one means the
  network never left its linearised regime. Panel~B shows the disorder
  susceptibility, whose interior maximum is present at every width.%
}{fig:lazy-rich}

\paragraph{Observation.}
Sweeping the initialisation scale over
$\AtlasLazyRichControlRange{}$ at \AtlasLazyRichSizes{} widths, the
initial-to-final kernel alignment rises from roughly $0.5$ to $0.99$ while the
relative weight movement \eqref{eq:lazy} falls by more than an order of
magnitude over the same range. The order parameter spans
$\AtlasLazyRichSemanticRatio{}$ times its own $95\%$ interval --- the largest
signal-to-noise ratio in the release --- and the susceptibility has an interior
peak at \emph{every} width, with height growing as
$N^{\AtlasLazyRichSharpening{}}$. Verdict
\emph{\AtlasLazyRichVerdict{}}, grade \AtlasLazyRichGrade{}.

\paragraph{Interpretation.}
The lazy and rich regimes are separated by a genuinely critical response, not a
gentle crossover: the disorder susceptibility of the kernel-alignment order
parameter diverges with width along the sampled path. Both halves of the
operational definition move together, which is the check that the alignment is
measuring feature learning rather than an optimisation artefact.

\paragraph{Not established.}
The location. The peak drifts and its extrapolation is inadmissible, so we
report a critical response without naming a critical initialisation scale.
Nothing here identifies the exponent with a known universality class.

\paragraph{Next falsifier.}
Extend the width path by a further factor of four. If the peak positions
stabilise, the drift fit becomes admissible and the crossover acquires a
location; if the peak height stops growing, the divergence is a small-width
effect and the verdict should fall to a crossover.

\subsection{Grokking as a delayed-generalisation transition}

\paragraph{Observation.}
An overparameterised student on a small training set fits almost immediately
and generalises much later, if at all. Across
\AtlasGrokkingConditions{} conditions the measured delay ranges over
$\AtlasGrokkingSemanticRatio{}$ times its interval, and the susceptibility of
the delay has an interior peak at every training-set size with height growing as
$N^{\AtlasGrokkingSharpening{}}$. Verdict \emph{\AtlasGrokkingVerdict{}}.

\paragraph{Interpretation.}
Delayed generalisation behaves like a critical phenomenon in the weight-decay
control: the run-to-run spread of the delay --- not merely its mean --- peaks in
the interior and grows with training-set size. That is a stronger statement than
the usual demonstration that a delay exists, and it is the quantity a theory of
grokking would have to reproduce.

\paragraph{Not established.}
A critical weight decay. Much of the grid is censored by construction: below a
threshold the student never generalises within the budget and above it the
delay vanishes, so the delay is only defined in a window. Censored runs are
flagged rather than being assigned a delay of zero --- differencing two censored
endpoints would report the \emph{opposite} of what happened, and this was a real
defect in an earlier version of the estimator.

\paragraph{Next falsifier.}
Lengthen the training budget by an order of magnitude. If the censored region
shrinks and the peak stays put, the boundary is real; if the peak simply tracks
the budget, the delay is a property of the stopping rule.

\subsection{Double descent}

\paragraph{Observation.}
Sweeping the parameterisation ratio across the interpolation threshold at three
label-noise levels reproduces the non-monotone risk curve, with the order
parameter spanning $\AtlasDoubleDescentSemanticRatio{}$ times its interval.
Verdict \emph{\AtlasDoubleDescentVerdict{}}.

\paragraph{Interpretation.}
The peak in test risk is located on the control axis and grows with label
noise, as expected. But the second axis here is the noise level, which changes
the problem rather than the size of one system, so under the rule of
\S\ref{sec:gates} this family cannot make a finite-size claim at all. The
result is a located threshold with no critical content --- which is precisely
what most double-descent figures in the literature actually establish.

\paragraph{Not established.}
That the interpolation peak sharpens. Demonstrating that requires a genuine
width or sample-size scale path at fixed noise, which this design does not have.

\subsection{SGD noise as an effective temperature}

\paragraph{Observation.}
The control is the ratio $T\simeq\eta/B$ of learning rate to batch size, swept
so that several distinct $(\eta,B)$ pairs realise the same $T$. The stationary
risk varies over $\AtlasSgdTemperatureSemanticRatio{}$ times its interval and
its susceptibility peaks in the interior at every size, with height growing as
$N^{\AtlasSgdTemperatureSharpening{}}$. Verdict
\emph{\AtlasSgdTemperatureVerdict{}}, grade \AtlasSgdTemperatureGrade{}.

\paragraph{Interpretation.}
The stationary state responds critically to the effective temperature, and the
design permits the sharper question --- whether the state depends on
$\eta$ and $B$ only through their ratio --- because the same $T$ is realised by
several pairs. That collapse is the operational content of calling $\eta/B$ a
temperature at all.

\paragraph{Not established.}
A critical temperature, for the usual reason: the peak location does not
extrapolate within the sampled window. We also do not claim the
$(\eta,B)$ collapse is exact, only that the design can test it.

\subsection{Mode connectivity: a negative result}

\paragraph{Observation.}
Two students trained from independent initialisations on independent data are
interpolated linearly, and the normalised barrier between them is measured
against width at three depths. The barrier varies strongly with \emph{depth} ---
roughly $0.4$ at depth two and above $1.2$ at depth four --- but its variation
with width spans only $\AtlasModeConnectivitySemanticRatio{}$ times its $95\%$
interval. The verdict is \emph{\AtlasModeConnectivityVerdict{}}, grade
\AtlasModeConnectivityGrade{}.

\paragraph{Interpretation.}
At these widths and this replication, the barrier's width dependence does not
separate from seed noise. We therefore report no width trend, in either
direction. The depth dependence is large and unambiguous, but depth is the
scale axis of this design, so it is a between-curve comparison rather than a
finite-size trend.

\paragraph{Not established.}
Linear mode connectivity, in either direction. In particular this is
\emph{not} evidence that barriers fail to vanish with width; it is evidence
that the present design cannot see the trend.

\paragraph{Next falsifier.}
Twelve seeds and a width path extended upward by a factor of four, at fixed
depth two, where the measured trend is at least in the expected direction.
